% Generated from paper/full_draft. Edit the Markdown source or migration script.

\section{Results}
\label{sec:06-results:results}

\subsection{Main Calibrated Token-Quality Frontier}
\label{sec:06-results:main-calibrated-token-quality-frontier}

The main cost result uses a calibrated token-budget policy. For each corpus scale, the budget is selected on calibration queries and then frozen before test evaluation. The measured cost is final LLM evidence-context input tokens relative to dense top-10, not retrieval-side candidate count.

\begin{table*}[t]
\centering
\small
\caption{Calibrated token-quality frontier on LoTTE technology/search}
\label{tab:2}
\resizebox{\textwidth}{!}{%
\begin{tabular}{lllrlr}
\toprule
Scale & Frozen budget policy & IntentWeight hit delta vs dense & IntentWeight token saving & Dense adaptive hit delta & Dense adaptive token saving \\
\midrule
100k & \nolinkurl{token_budget_r0.95_m4} & +0.00 pp & 6.18\% & -1.44 pp & 13.83\% \\
200k & \nolinkurl{token_budget_r0.85_m4} & +1.20 pp & 16.00\% & -2.40 pp & 21.95\% \\
400k & \nolinkurl{token_budget_r0.98_m4} & +2.32 pp & 6.57\% & -0.24 pp & 11.44\% \\
638k & \nolinkurl{token_budget_r0.85_m4} & -0.08 pp & 17.53\% & -3.84 pp & 21.90\% \\
\bottomrule
\end{tabular}%
}
\end{table*}

The calibrated IntentWeight policies save 6-18\% final evidence-context tokens under frozen policy selection. Dense-only adaptive truncation usually saves more tokens, but it loses $\mathrm{Hit@10}$ on every scale. This comparison shows that the effect is not merely dense top-k truncation: route confidence helps decide where a shorter context remains safe. Strict seed-level non-inferiority remains scale-dependent, so the paper should claim a bounded quality-cost frontier rather than universal statistical superiority.

The earlier conservative confidence-only policy remains useful as a stable baseline. It reduces final retrieved context tokens by 4.7-5.3\% across LoTTE technology/search 100k-638k while preserving dense-level $\mathrm{Hit@10}$. Appendix A reports the complete confidence-only scale table and seed stability diagnostics.

\subsection{Cross-Domain Validation}
\label{sec:06-results:cross-domain-validation}

We replicate the validation pattern on LoTTE science/search to test whether the effect is limited to technology/search.

\begin{table}[t]
\centering
\small
\caption{Cross-domain validation on LoTTE science/search}
\label{tab:3}
\resizebox{\linewidth}{!}{%
\begin{tabular}{lrrlr}
\toprule
Domain/scale & Dense $\mathrm{Hit@10}$ & IntentWeight fixed top-10 $\mathrm{Hit@10}$ & Hit delta & Budgeted token saving \\
\midrule
science/search 20k/q200 & 0.8950 & 0.9267 & +3.17 pp & 13.18-14.31\% \\
science/search 100k & 0.8926 & 0.9077 & +1.51 pp & 17.53-20.53\% \\
\bottomrule
\end{tabular}%
}
\end{table}

The fixed top-10 ranking-side effect transfers to a second LoTTE domain. The context-budget result is more nuanced. At 20k/q200, the frozen budgeted policies keep above-dense $\mathrm{Hit@10}$ while saving roughly 13-14\% context tokens. At 100k, the same aggressive budget saves roughly 17-21\% tokens but can introduce small $\mathrm{Hit@10}$ drops. This supports adaptive ranking across domains while showing that compression strength must be calibrated per domain and scale.

\subsection{Component Ablation}
\label{sec:06-results:component-ablation}

The component ablation summarizes which parts of the system provide the quality floor, routing signal, feedback adaptation, and final token saving on LoTTE 100k.

\begin{table*}[t]
\centering
\small
\caption{LoTTE 100k component ablation. The no-feedback gated row disables learning; its high $\mathrm{Hit@10}$ reflects full dense fallback ($\mathrm{dense\ rate}=1.0$), not learned route efficiency}
\label{tab:4}
\resizebox{\textwidth}{!}{%
\begin{tabular}{llrrrrrrrr}
\toprule
Component & Role & $\mathrm{Hit@10}$ & $\mathrm{EvidenceRecall@10}$ & $\mathrm{Tokens@10}$ & Token ratio & Dense rate & LinUCB rate & Cluster hit & Last reward \\
\midrule
Dense-only & Quality floor & 0.8674 & 0.7026 & 1472.39 & 1.0000 & - & - & - & - \\
BM25-only & Lexical baseline & 0.7232 & 0.5240 & 1745.12 & 1.1852 & - & - & - & - \\
Dense+BM25 hybrid & Static fusion & 0.8624 & 0.6848 & 1705.46 & 1.1583 & - & - & - & - \\
No feedback gated & Full dense fallback, no learning & 0.8826 & 0.7246 & 1561.15 & 1.0603 & 1.0000 & 0.0000 & 0.1553 & 0.1516 \\
Equal noisy feedback & No trust weighting & 0.8641 & 0.6604 & 1423.84 & 0.9670 & 0.7480 & 0.2520 & 0.5979 & 0.7517 \\
Trust-weighted feedback & Default trust scoring & 0.8641 & 0.6661 & 1399.51 & 0.9505 & 0.6708 & 0.3292 & 0.7223 & 0.8328 \\
Trust-weighted mild noise & Best controlled-noise point & 0.8775 & 0.6795 & 1362.68 & 0.9255 & 0.5826 & 0.4174 & 0.7908 & 0.8820 \\
Conservative final policy & Confidence-only baseline & 0.8652 & 0.6737 & 1401.24 & 0.9517 & 0.6708 & 0.3292 & 0.7223 & 0.8328 \\
Oracle feedback & Upper bound & 0.8758 & 0.6768 & 1327.03 & 0.9013 & 0.4345 & 0.5655 & 0.8386 & 0.8932 \\
\bottomrule
\end{tabular}%
}
\end{table*}

Dense-only remains the quality floor. BM25-only is weaker as a standalone retriever, and static dense+BM25 hybrid is near dense but uses more final context tokens. Trust-weighted feedback improves policy internals relative to equal noisy feedback, especially selected-cluster hit and last true reward. The oracle row shows the upper bound under clean feedback.

\subsection{Feedback-Driven Adaptation and Recovery}
\label{sec:06-results:feedback-driven-adaptation-and-recovery}

Final $\mathrm{Hit@10}$ can be saturated by dense and BM25 rescue routes, making feedback gains less visible in the fused final ranking. The clearer feedback signal appears in route-policy metrics such as selected-cluster hit, last true reward, dense rate, and LinUCB usage.

\begin{table*}[t]
\centering
\small
\caption{Feedback-driven policy adaptation on LoTTE 100k}
\label{tab:5}
\resizebox{\textwidth}{!}{%
\begin{tabular}{lrrrrrr}
\toprule
Feedback mode & $\mathrm{Hit@10}$ & Token ratio & Dense rate & LinUCB rate & Selected-cluster hit & Last true reward \\
\midrule
Equal noisy feedback & 0.8641 & 0.9670 & 0.7480 & 0.2520 & 0.5979 & 0.7517 \\
Trust-weighted feedback & 0.8641 & 0.9505 & 0.6708 & 0.3292 & 0.7223 & 0.8328 \\
Trust-weighted mild noise & 0.8775 & 0.9255 & 0.5826 & 0.4174 & 0.7908 & 0.8820 \\
Oracle feedback & 0.8758 & 0.9013 & 0.4345 & 0.5655 & 0.8386 & 0.8932 \\
\bottomrule
\end{tabular}%
}
\end{table*}

Under default noisy feedback, trust weighting improves selected-cluster hit from $0.5979$ to $0.7223$ and last true reward from $0.7517$ to $0.8328$ relative to equal noisy feedback. Under mild trust-weighted noise, selected-cluster hit reaches $0.7908$, last true reward reaches $0.8820$, dense rate falls to $0.5826$, and final context token ratio falls to $0.9255\times$.

Feedback also provides a recovery path for tail queries harmed by aggressive context budgets. A query is affected when dense top-10 retrieves at least one GT evidence chunk but the budgeted IntentWeight context misses. Same-query retry is a post-feedback repair setting, not a first-pass IID ranking claim.

\begin{table}[t]
\centering
\small
\caption{Conservative post-feedback recovery on affected LoTTE 100k queries}
\label{tab:6}
\resizebox{\linewidth}{!}{%
\begin{tabular}{lrrrr}
\toprule
Domain & Affected queries & Recovered & Recovery rate & Avg token saving vs dense \\
\midrule
science/search 100k & 34 & 14 & 41.18\% & 5.76\% \\
technology/search 100k & 42 & 9 & 21.43\% & 11.75\% \\
pooled & 76 & 23 & 30.26\% & - \\
\bottomrule
\end{tabular}%
}
\end{table}

The pooled conservative retry recovery rate is approximately 30\%, with an approximate Wilson interval around 21-41\%. This supports a practical recovery claim: budget-induced failures are not always permanent, and feedback can repair a meaningful fraction of them. A stricter calibration-to-test variant has only small, domain-dependent effects, so feedback should be framed as a controlled fallback trigger rather than as unconditional global reranking.

\subsection{Geometry Diagnostics}
\label{sec:06-results:geometry-diagnostics}

The geometry scale diagnostic validates whether LoTTE retains usable local geometry as scale grows.

\begin{table*}[t]
\centering
\small
\caption{LoTTE geometry diagnostics across corpus scale}
\label{tab:7}
\resizebox{\textwidth}{!}{%
\begin{tabular}{lrrrrl}
\toprule
Scale & $\mathrm{PCAdim90}$ sample & $\mathrm{PCAvar@64}$ sample & $\mathrm{NearestClusterHit@3}$ & $\mathrm{ContextRetention@10}$ & Confidence-only hit delta \\
\midrule
100k & 182 & 0.6437 & 0.8870 & 0.9033 & -0.22 pp \\
200k & 186 & 0.6292 & 0.8697 & 0.8947 & +2.80 pp \\
400k & 190 & 0.6110 & 0.9016 & 0.8826 & +1.01 pp \\
638k & 196 & 0.5867 & 0.9016 & 0.8571 & +1.85 pp \\
\bottomrule
\end{tabular}%
}
\end{table*}

$\mathrm{NearestClusterHit@3}$ remains high, around 0.87-0.90, suggesting local geometry is useful for routing. $\mathrm{PCAdim90}$ increases and $\mathrm{PCAvar@64}$ decreases with scale, suggesting the representation geometry becomes more complex. Context retention declines with scale, showing that geometry alone should not replace dense retrieval.

These diagnostics support the piecewise relevance-manifold framing as a useful motivation and diagnostic, not as a theorem. Figure~\ref{fig:geometry} visualizes the same trend.

\subsection{Boundary, Robustness, and Downstream Checks}
\label{sec:06-results:boundary-robustness-and-downstream-checks}

Secondary datasets and checks are used to bound the claim rather than expand it into universal dense-retrieval dominance.

PubMedQA and Banking77 support the feedback-adaptation mechanism near quality ceilings. PubMedQA is an evidence-retrieval proof-of-concept with section-level ground truth, while Banking77 is better understood as an intent-routing proxy. eManual and CUAD are limitation cases: eManual has heavy duplicate evidence text and strict chunk-id labels, while CUAD remains a sparse GT-anchored smoke case.

The encoder robustness check uses \nolinkurl{sentence-transformers/multi-qa-MiniLM-L6-cos-v1}. Under this stronger MiniLM-family encoder, the dense baseline improves to $\mathrm{Hit@10}=0.8809$, and IntentWeight still reaches mean $\mathrm{Hit@10}=0.8853$ while reducing final context tokens by 3.35\%. The downstream answer-quality check compares 60 LoTTE 100k queries with dense top-10 and compressed context; it does not show obvious answer-quality degradation, but it remains a small sanity check rather than a full human evaluation. Appendix D-F report the full boundary, robustness, and downstream tables.
